\chapter{Materiales y degradación de aeromotores}\label{ch:10}
\objetivos{Explicar por qué cambian los materiales a lo largo del motor y relacionar cada pieza con sus solicitaciones y mecanismos de daño.}
\section{Un mapa de temperatura y función}
El recorrido del aire y del gas atraviesa fan, compresores, cámara de combustión, turbinas y salida. La temperatura del gas aumenta con la compresión y, especialmente, con la combustión. La temperatura del metal depende además de refrigeración, recubrimientos, transferencia térmica y régimen transitorio. No se iguala automáticamente a una temperatura medida del gas.

La selección se realiza por componente. Una carcasa puede priorizar contención y estabilidad dimensional; un disco, integridad bajo esfuerzos centrífugos y ciclos; un álabe, aerodinámica, vibración y resistencia térmica; un cojinete, contacto, desgaste y lubricación. La exposición al daño por objetos extraños y el acceso de inspección añaden restricciones.
\ilustra{motor-materiales}{Mapa funcional de un turbofán y de familias de materiales que pueden emplearse en distintas zonas. Esquema de tendencias, no despiece de un motor concreto.}
\section{Fan y compresor}
El fan debe combinar masa, rigidez, resistencia centrífuga, vibración e impacto. Según diseño se utilizan aleaciones de titanio, compuestos con protecciones y otras soluciones. El compresor añade incrementos de temperatura, esfuerzos cíclicos y riesgos asociados a rozamientos o ingestión. Las últimas etapas pueden exigir familias de mayor capacidad térmica.

Erosión y depósitos alteran perfiles, holguras y rugosidad; con ello pueden modificar rendimiento y margen de operación. Una pérdida de eficiencia no identifica por sí sola una grieta. Una entalla por impacto puede ser pequeña en masa perdida y relevante para fatiga, por lo que las dimensiones y ubicación del daño importan.
\section{Cámara de combustión y turbina}
La cámara combina llama, gradientes, oxidación, vibración y fatiga térmica. El revestimiento metálico y sus recubrimientos trabajan junto a refrigeración y geometría de orificios. Un deterioro local de refrigeración puede elevar temperatura del metal aunque una medida global siga dentro de una banda habitual.

Los álabes de turbina pueden experimentar fluencia y fatiga termomecánica, además de oxidación, corrosión caliente y erosión. La solidificación direccional orienta granos para reducir ciertos límites transversales; los monocristales eliminan límites de alto ángulo en el cristal, pero exigen control de orientación, porosidad y defectos de colada. No eliminan toda posibilidad de daño ni la necesidad de inspección.
\section{Discos, carcasas y elementos de contacto}
La integridad de un disco necesita gestionar esfuerzos centrífugos, gradientes y ciclos. La resistencia a iniciación y propagación de defectos debe evaluarse junto a la dispersión de fabricación. En un disco de metal procesado por polvo, la limpieza, las inclusiones y la homogeneidad de proceso pueden ser críticas; la composición nominal es insuficiente.

Las carcasas condicionan holguras y alineación. Su dilatación respecto a rotores puede variar durante arranques y paradas, dando lugar a contactos. En cojinetes, la película lubricante, el estado de superficies, la contaminación y la temperatura se relacionan con desgaste, fatiga de contacto y calentamiento. Un análisis de partículas en aceite aporta evidencia de desgaste, pero necesita origen y tendencia.
\section{Barreras térmicas y ambientales}
Una barrera térmica suele combinar sustrato, capa de enlace y capa cerámica con una arquitectura determinada. La diferencia de expansión térmica, el crecimiento de óxidos en la interfaz, el ciclado y los depósitos pueden favorecer desprendimientos. El recubrimiento reduce solicitación térmica bajo condiciones de transferencia determinadas; no garantiza una reducción fija de temperatura en todos los motores.

En ciertos CMC, los recubrimientos ambientales protegen frente a mecanismos químicos, por ejemplo asociados al ambiente de combustión. No debe confundirse protección ambiental con aislamiento térmico, aunque puedan coexistir funciones. La evaluación de una zona sin recubrimiento necesita datos específicos de ubicación y extensión.
\caso{M10. Una subida de temperatura no diagnostica sola}{En un motor ficticio aumenta la temperatura de escape necesaria para mantener un empuje comparable. Entre las hipótesis figuran deterioro de perfiles, mayores holguras, suciedad, cambio de ambiente o error de sensor. Se comparan parámetros corregidos y condiciones de misión antes de atribuirlo a una turbina degradada. Si además aparece un hallazgo de boroscopia, se relaciona su ubicación con la tendencia. El caso enseña a combinar evidencias sin inferir un mecanismo único de una sola señal.}
\section{Del mecanismo al plan de seguimiento}
Un mapa de daño vincula pieza, mecanismo, factor impulsor, indicio y técnica. Para fatiga de un disco puede interesar iniciación en una zona concentradora y un método sensible al defecto esperado; para un revestimiento caliente, la distribución de grietas o pérdida de recubrimiento puede ser la evidencia principal. Los límites reales dependen de los datos de mantenimiento del motor.

La ingeniería de materiales y el mantenimiento se encuentran en la trazabilidad: qué pieza, qué estado, qué historial y qué datos justifican una acción. La modelización académica ayuda a interpretar por qué existe un límite o una inspección; no modifica ese límite ni sustituye la instrucción aplicable.

\section{Un motor impone campos, no una sola temperatura}
La temperatura del gas, la del metal y la de un sensor no son intercambiables. La refrigeración interna y superficial, el espesor, la conductividad y los recubrimientos crean gradientes. Un aumento de temperatura de gas no permite sumar el mismo incremento a todos los componentes. El cálculo térmico también necesita condiciones de contorno, transitorios y radiación cuando sea relevante.

Los gradientes generan deformaciones diferenciales; las restricciones convierten parte de ellas en tensión. En un arranque o parada pueden coexistir tensiones mecánicas y térmicas de signos y distribuciones cambiantes. La compatibilidad entre disco y álabe, holguras con carcasa y evolución del recubrimiento relaciona materiales y diseño térmico. Optimizar una propiedad aislada puede trasladar el problema a otra interfaz.
\section{Microestructura orientada, refrigeración y recubrimiento}
La solidificación direccional reduce ciertas fronteras transversales a la solicitación; un monocristal evita bordes de grano de alto ángulo dentro de la idealización del producto, pero conserva anisotropía cristalina y otros defectos posibles. Orientación, defectos de colada, tratamiento y acabado necesitan control. No equivale a un material sin mecanismos de fallo.

Un sistema de barrera térmica incluye funciones del recubrimiento cerámico, capa de unión y óxido que evoluciona durante exposición, según arquitectura. Diferencias de dilatación, oxidación, erosión y ciclos pueden iniciar pérdida de adherencia. Un recubrimiento desprendido modifica el campo térmico del sustrato; medir solo su espesor medio no describe todas las condiciones de protección. El ejemplo térmico P35 enseña resistencias en serie, sin reproducir una geometría refrigerada real.
\section{Un mapa de daño por componente}
En fan y compresor interesan impacto, erosión, fatiga y degradación aerodinámica; el material y la etapa condicionan la relevancia de cada uno. En cámaras y zonas calientes aparecen ciclos térmicos, oxidación, corrosión caliente, fisuración y degradación de recubrimientos. En discos se estudian fatiga y defectos en un campo centrífugo y térmico. En cojinetes, contacto, lubricación y contaminación exigen otro enfoque.

Este mapa se utiliza para formular preguntas de inspección: dónde, qué mecanismo, qué evidencia y qué método lo puede observar. No convierte una lista de daños posibles en diagnóstico. Dos mecanismos pueden producir señales parecidas; por ejemplo, una pérdida de rendimiento puede deberse a degradación de perfiles, holguras o condiciones de medida. El capítulo~\ref{ch:22} conecta este razonamiento con seguimiento en servicio.
\repaso{¿Por qué una cerámica puede trabajar en una zona caliente y aun así requerir protección ambiental?}{La estabilidad térmica no garantiza resistencia frente a todas las especies del gas, vapor, depósitos, erosión o ciclos. La selección considera química del ambiente y sistema completo de protección, además de la temperatura.}

\resumen{La temperatura del gas y la del metal son distintas. Fan, compresor, cámara, turbina y discos no comparten un criterio único de selección. Monocristales y recubrimientos controlan mecanismos concretos; no suprimen todos los modos de fallo. La monitorización necesita contexto y corroboración.}
\fuentes{NASA, superaleaciones y materiales cerámicos \cite{NASA-SX,NASA-CMC}; FAA--General, capítulos 7, 8 y 10 \cite{FAA-G}; documentos docentes de materiales.}
